% GENERATED FILE. Edit canonical lessons, then run npm run generate:books.
% canonical-source-hash: fnv1a64:69ee9a4a364a4221
% canonical-lessons: RU-C221-igrushka, RU-C221-myach, RU-C221-kukla, RU-C221-vaza, RU-C221-kartina, RU-R221-first-pass-words-from-chapters-213-216, RU-R221-second-pass-words-from-chapters-217-221

\chapter{Toy, Ball, Doll, Vase, Painting}
\label{ch:ru-toy-ball-doll-vase-painting}
\hlchaptermodality{221}

\begin{chapteropening}
\textbf{By the end of this chapter:} \emph{I can say a toy, a ball, a doll, a vase and a painting.}

Complete the last lesson of chapter 221: I can say a toy, a ball, a doll, a vase and a painting.
\end{chapteropening}

\section[\texorpdfstring{\ru{игрушка} (igrúshka)}{игрушка (igrúshka)}]{\ru{игрушка} (igrúshka) — a toy}

\label{lesson:RU-C221-igrushka}



\begin{quote}
Before the new one: say the Russian for a printer, then the Russian for a remote control.
\end{quote}

\subsection*{You'll want to know: \ru{игрушка}}
\textbf{\ru{игрушка}} (\emph{igrúshka}) — \textquotedblleft{}a toy\textquotedblright{}.

\begin{grammarlens}[title={What you've built}]
One more word for this chapter.
\end{grammarlens}

\subsection*{Your turn}
\begin{itemize}
\raggedright
  \item \emph{Say it:} \emph{igrúshka}
  \item \emph{Say it:} \emph{igrúshka}, once more
  \item \emph{Say it:} say \emph{igrúshka}
  \item \emph{Recall:} say the Russian for a printer, then the Russian for a remote control, then say \emph{igrúshka} again
\end{itemize}

\begin{tcolorbox}[breakable,colback=teal!4,colframe=teal!35!black,title={Before you move on}]
What does \ru{игрушка} mean? (\textquotedblleft{}A toy\textquotedblright{}.) Say it once more.
\end{tcolorbox}
\section[\texorpdfstring{\ru{мяч} (myach)}{мяч (myach)}]{\ru{мяч} (myach) — a ball}

\label{lesson:RU-C221-myach}



\begin{quote}
Before the new one: say the Russian for a remote control, then the Russian for a toy.
\end{quote}

\subsection*{You'll want to know: \ru{мяч}}
\textbf{\ru{мяч}} (\emph{myach}) — \textquotedblleft{}a ball\textquotedblright{}.

\begin{grammarlens}[title={What you've built}]
One more word for this chapter.
\end{grammarlens}

\subsection*{Your turn}
\begin{itemize}
\raggedright
  \item \emph{Say it:} \emph{myach}
  \item \emph{Say it:} \emph{myach}, once more
  \item \emph{Say it:} say \emph{myach}
  \item \emph{Recall:} say the Russian for a remote control, then the Russian for a toy, then say \emph{myach} again
\end{itemize}

\begin{tcolorbox}[breakable,colback=teal!4,colframe=teal!35!black,title={Before you move on}]
What does \ru{мяч} mean? (\textquotedblleft{}A ball\textquotedblright{}.) Say it once more.
\end{tcolorbox}
\section[\texorpdfstring{\ru{кукла} (kúkla)}{кукла (kúkla)}]{\ru{кукла} (kúkla) — a doll}

\label{lesson:RU-C221-kukla}



\begin{quote}
Before the new one: say the Russian for a toy, then the Russian for a ball.
\end{quote}

\subsection*{You'll want to know: \ru{кукла}}
\textbf{\ru{кукла}} (\emph{kúkla}) — \textquotedblleft{}a doll\textquotedblright{}.

\begin{grammarlens}[title={What you've built}]
One more word for this chapter.
\end{grammarlens}

\subsection*{Your turn}
\begin{itemize}
\raggedright
  \item \emph{Say it:} \emph{kúkla}
  \item \emph{Say it:} \emph{kúkla}, once more
  \item \emph{Say it:} say \emph{kúkla}
  \item \emph{Recall:} say the Russian for a toy, then the Russian for a ball, then say \emph{kúkla} again
\end{itemize}

\begin{tcolorbox}[breakable,colback=teal!4,colframe=teal!35!black,title={Before you move on}]
What does \ru{кукла} mean? (\textquotedblleft{}A doll\textquotedblright{}.) Say it once more.
\end{tcolorbox}
\section[\texorpdfstring{\ru{ваза} (váza)}{ваза (váza)}]{\ru{ваза} (váza) — a vase}

\label{lesson:RU-C221-vaza}



\begin{quote}
Before the new one: say the Russian for a ball, then the Russian for a doll.
\end{quote}

\subsection*{You'll want to know: \ru{ваза}}
\textbf{\ru{ваза}} (\emph{váza}) — \textquotedblleft{}a vase\textquotedblright{}.

\begin{grammarlens}[title={What you've built}]
One more word for this chapter.
\end{grammarlens}

\subsection*{Your turn}
\begin{itemize}
\raggedright
  \item \emph{Say it:} \emph{váza}
  \item \emph{Say it:} \emph{váza}, once more
  \item \emph{Say it:} say \emph{váza}
  \item \emph{Recall:} say the Russian for a ball, then the Russian for a doll, then say \emph{váza} again
\end{itemize}

\begin{tcolorbox}[breakable,colback=teal!4,colframe=teal!35!black,title={Before you move on}]
What does \ru{ваза} mean? (\textquotedblleft{}A vase\textquotedblright{}.) Say it once more.
\end{tcolorbox}
\section[\texorpdfstring{\ru{картина} (kartína)}{картина (kartína)}]{\ru{картина} (kartína) — a painting}

\label{lesson:RU-C221-kartina}



\begin{quote}
Before the new one: say the Russian for a doll, then the Russian for a vase.
\end{quote}

\subsection*{You'll want to know: \ru{картина}}
\textbf{\ru{картина}} (\emph{kartína}) — \textquotedblleft{}a painting\textquotedblright{}.

\begin{grammarlens}[title={What you've built}]
One more word for this chapter.
\end{grammarlens}

\subsection*{Your turn}
\begin{itemize}
\raggedright
  \item \emph{Say it:} \emph{kartína}
  \item \emph{Say it:} \emph{kartína}, once more
  \item \emph{Say it:} say \emph{kartína}
  \item \emph{Recall:} say the Russian for a doll, then the Russian for a vase, then say \emph{kartína} again
\end{itemize}

\begin{tcolorbox}[breakable,colback=teal!4,colframe=teal!35!black,title={Before you move on}]
What does \ru{картина} mean? (\textquotedblleft{}A painting\textquotedblright{}.) Say it once more.
\end{tcolorbox}
\section[(dialogue)]{Words from chapters 213-216 — a first pass}

\label{lesson:RU-R221-first-pass-words-from-chapters-213-216}



\begin{quote}
Say the words for a vase and a painting.
\end{quote}

\subsection*{Your turn}
Say these aloud:

\begin{itemize}
\raggedright
  \item a border, customs, a capital, a flat, a floor
  \item from here, back, forward, up, down
  \item of course, probably, maybe, unfortunately, fortunately
  \item calm, nervous, happy, sad, strict
\end{itemize}

\begin{tcolorbox}[breakable,colback=teal!4,colframe=teal!35!black,title={Before you move on}]
A border? (\textbf{\ru{граница}}.) Customs? (\textbf{\ru{таможня}}.) A capital? (\textbf{\ru{столица}}.) A flat? (\textbf{\ru{квартира}}.) A floor? (\textbf{\ru{этаж}}.) From here? (\textbf{\ru{отсюда}}.) Back? (\textbf{\ru{назад}}.) Forward? (\textbf{\ru{вперёд}}.) Up? (\textbf{\ru{вверх}}.) Down? (\textbf{\ru{вниз}}.) Of course? (\textbf{\ru{конечно}}.) Probably? (\textbf{\ru{наверное}}.) Maybe? (\textbf{\ru{может} \ru{быть}}.) Unfortunately? (\textbf{\ru{к} \ru{сожалению}}.) Fortunately? (\textbf{\ru{к} \ru{счастью}}.) Calm? (\textbf{\ru{спокойный}}.) Nervous? (\textbf{\ru{нервный}}.) Happy? (\textbf{\ru{счастливый}}.) Sad? (\textbf{\ru{грустный}}.) Strict? (\textbf{\ru{строгий}}.)
\end{tcolorbox}
\section[(dialogue)]{Words from chapters 217-221 — a second pass}

\label{lesson:RU-R221-second-pass-words-from-chapters-217-221}



\begin{quote}
Say the words for blunt and smooth.
\end{quote}

\subsection*{Your turn}
Say these aloud:

\begin{itemize}
\raggedright
  \item blunt, smooth, cool, fresh, modern
  \item different, the same, similar, main, ordinary
  \item toothpaste, a razor, an iron, a vacuum cleaner, an oven
  \item a tablet, a screen, a keyboard, a printer, a remote control
  \item a toy, a ball, a doll, a vase, a painting
\end{itemize}

\begin{tcolorbox}[breakable,colback=teal!4,colframe=teal!35!black,title={Before you move on}]
Blunt? (\textbf{\ru{тупой}}.) Smooth? (\textbf{\ru{гладкий}}.) Cool? (\textbf{\ru{прохладный}}.) Fresh? (\textbf{\ru{свежий}}.) Modern? (\textbf{\ru{современный}}.) Different? (\textbf{\ru{разный}}.) The same? (\textbf{\ru{одинаковый}}.) Similar? (\textbf{\ru{похожий}}.) Main? (\textbf{\ru{главный}}.) Ordinary? (\textbf{\ru{обычный}}.) Toothpaste? (\textbf{\ru{зубная} \ru{паста}}.) A razor? (\textbf{\ru{бритва}}.) An iron? (\textbf{\ru{утюг}}.) A vacuum cleaner? (\textbf{\ru{пылесос}}.) An oven? (\textbf{\ru{духовка}}.) A tablet? (\textbf{\ru{планшет}}.) A screen? (\textbf{\ru{экран}}.) A keyboard? (\textbf{\ru{клавиатура}}.) A printer? (\textbf{\ru{принтер}}.) A remote control? (\textbf{\ru{пульт}}.) A toy? (\textbf{\ru{игрушка}}.) A ball? (\textbf{\ru{мяч}}.) A doll? (\textbf{\ru{кукла}}.) A vase? (\textbf{\ru{ваза}}.) A painting? (\textbf{\ru{картина}}.)
\end{tcolorbox}
